\chapter{Introdução}
\label{chap:introducao}

Este capítulo contém um texto genérico para introduzir algumas referências básicas de Física Médica (e assim ter um exemplo de Bibliografia no final do documento).

\section{Contexto}
\label{sec:contexto}

A Física Médica combina fundamentos de interação da radiação com a matéria, dosimetria, instrumentação, formação de imagens, radiobiologia, proteção radiológica e aplicação clínica. Para os fundamentos de física radiológica e dosimetria, uma referência clássica é o texto de Attix, que desenvolve grandezas dosimétricas, equilíbrio de partículas carregadas, teoria de cavidade e detectores empregados em dosimetria\supercite{attix1986}. O manual de física de radioterapia editado por Podgorsak oferece uma visão ampla da física das radiações e de sua aplicação clínica\supercite{podgorsak2005}. Em espanhol, a coleção \emph{Fundamentos de Física Médica}, editada pela Sociedad Española de Física Médica, constitui uma referência estruturada para a formação na área\supercite{brosed2011}. Para instrumentação e detecção, o livro de Knoll permanece uma das referências centrais sobre mecanismos de detecção, estatística de contagem e espectrometria\supercite{knoll2010}, enquanto Cember e Johnson apresentam os fundamentos de \emph{health physics} e proteção radiológica aplicada\supercite{cember2008}.

\section{Radiodiagnóstico e formação de imagens}

Para física do radiodiagnóstico, \emph{The Essential Physics of Medical Imaging} reúne fundamentos de radiografia, fluoroscopia, mamografia, tomografia computadorizada, medicina nuclear, ressonância magnética e ultrassonografia\supercite{bushberg2020}. O handbook da IAEA para física do radiodiagnóstico fornece uma abordagem sistemática para princípios físicos, dosimetria e controle de qualidade\supercite{iaea_diagnostic2014}. Hendee e Ritenour apresentam uma abordagem clássica da física da formação de imagens médicas e de suas diferentes modalidades\supercite{hendee2002}. Para aprofundamento em detectores semicondutores e eletrônica de leitura, o texto de Spieler é particularmente útil\supercite{spieler2005}.

\section{Tomografia computadorizada}

Em tomografia computadorizada, Seeram apresenta os princípios físicos da aquisição, reconstrução, aplicações clínicas e controle de qualidade\supercite{seeram2015}. O AAPM TG-111 introduz uma metodologia abrangente para avaliação de dose em tomografia computadorizada, incluindo aquisições axiais, helicoidais e de feixe cônico\supercite{aapm_tg111_2010}. Para avaliação de desempenho dos sistemas modernos de TC, o AAPM TG-233 enfatiza métricas de qualidade de imagem dependentes da tarefa e metodologias apropriadas para reconstruções iterativas\supercite{aapm_tg233_2019}. A ICRP Publication 102 discute especificamente o gerenciamento da dose do paciente em tomografia computadorizada multidetectores\supercite{icrp102_2007}.

\section{Mamografia}

O programa de garantia da qualidade em mamografia digital da IAEA fornece procedimentos e critérios para avaliação sistemática do desempenho de equipamentos\supercite{iaea_mammo2011}. O protocolo de mamografia da EFOMP propõe um conjunto harmonizado de testes de controle de qualidade para sistemas digitais\supercite{efomp_mammo2015}; sua metodologia e fundamentação também foram apresentadas em publicação específica do grupo de trabalho da EFOMP\supercite{gennaro2018}. O \emph{Protocolo Español de Control de Calidad en Radiodiagnóstico} reúne procedimentos de controle de qualidade em diversas modalidades e contém uma seção dedicada à mamografia\supercite{sefm_sepr_seram2011}. O manual de controle de qualidade em mamografia digital do American College of Radiology constitui outra referência importante para programas de garantia da qualidade\supercite{acr_mammo2018}.

\section{Ressonância magnética}

Para ressonância magnética, McRobbie e colaboradores oferecem uma introdução que conecta a física de spins, sequências de pulsos e mecanismos de contraste à formação da imagem\supercite{mcrobbie2017}. O texto de Brown e colaboradores apresenta um tratamento mais aprofundado dos princípios físicos e do projeto de sequências de RM\supercite{brown2014}. Para testes de aceitação e programas de garantia da qualidade em instalações de ressonância magnética, o AAPM Report 100 constitui uma referência específica\supercite{aapm100_2010}.

\section{Medicina nuclear}

Em medicina nuclear, o livro de Cherry, Sorenson e Phelps cobre física de radionuclídeos, detectores, gama-câmaras, SPECT, PET, reconstrução e dosimetria interna\supercite{cherry2012}. O handbook da IAEA complementa essa formação com uma abordagem voltada ao treinamento do físico médico em medicina nuclear\supercite{iaea_nucmed2015}. A série de publicações do Committee on Medical Internal Radiation Dose fornece métodos e ferramentas padronizados para dosimetria interna\supercite{snmmi_mird}; entre os documentos metodológicos centrais, o MIRD Pamphlet No. 21 formaliza a nomenclatura e o esquema geral de dosimetria de radiofármacos\supercite{bolch2009}. A ICRP Publication 107 reúne dados de decaimento nuclear destinados a cálculos dosimétricos\supercite{icrp107_2008}, enquanto a ICRP Publication 128 compila coeficientes de dose e modelos biocinéticos para radiofármacos de uso frequente em diagnóstico\supercite{icrp128_2015}.

\section{Radioterapia}

Na física clínica da radioterapia, a sexta edição de \emph{Khan's The Physics of Radiation Therapy} apresenta dosimetria, planejamento, aceleradores lineares e técnicas modernas de tratamento\supercite{gibbons2019}. O código de prática IAEA TRS-398, em sua revisão de 2024, é uma referência central para determinação de dose absorvida em feixes externos com calibração baseada em dose absorvida na água\supercite{iaea_trs398_2024}. Para campos pequenos, o IAEA--AAPM TRS-483 estabelece um formalismo e procedimentos específicos de dosimetria de referência e relativa\supercite{iaea_trs483_2017}. Nos Estados Unidos, o protocolo AAPM TG-51 estabelece a dosimetria clínica de referência para feixes de fótons e elétrons de alta energia\supercite{almond1999}, tendo recebido posteriormente um adendo específico para feixes de fótons\supercite{mcewen2014}.

No controle de qualidade de aceleradores médicos, o AAPM TG-142 estabeleceu uma estrutura de testes e tolerâncias para aceleradores e sistemas de imagem associados\supercite{klein2009}. Para a verificação dosimétrica de tratamentos IMRT, o AAPM TG-218 apresenta metodologias e limites de tolerância para QA baseado em medidas\supercite{miften2018}. Em radioterapia estereotáxica extracraniana, o AAPM TG-101 reúne recomendações para implementação e operação clínica de programas de SBRT\supercite{benedict2010}.

A padronização da prescrição, do registro e do relato em radioterapia é tratada por uma sequência de documentos da ICRU. O Report 50 estabeleceu os conceitos fundamentais de volumes e relato de dose em radioterapia com fótons\supercite{icru50_1993}; o Report 62 complementou essas recomendações para técnicas mais modernas\supercite{icru62_1999}; o Report 83 adaptou o formalismo para radioterapia de intensidade modulada\supercite{icru83_2010}; e o Report 91 tratou especificamente de tratamentos estereotáxicos com pequenos feixes de fótons\supercite{icru91_2014}.

\section{Radiobiologia}

O livro de Hall e Giaccia é uma referência abrangente para radiobiologia celular e molecular aplicada à radiologia e à radioterapia\supercite{hall2018}. Para uma abordagem diretamente orientada à prática clínica, \emph{Basic Clinical Radiobiology} apresenta os fundamentos de resposta tecidual, fracionamento, modelo linear-quadrático, hipóxia, repopulação e efeitos tardios\supercite{joiner2018}.

\section{Proteção radiológica}

A ICRP Publication 103 contém as recomendações gerais que estruturam o sistema contemporâneo de proteção radiológica\supercite{icrp103_2007}. A aplicação desses princípios às exposições médicas é desenvolvida de forma específica na ICRP Publication 105\supercite{icrp105_2007}. No âmbito dos padrões internacionais de segurança, o GSR Part 3 da IAEA estabelece os requisitos básicos de proteção radiológica e segurança de fontes de radiação\supercite{iaea_gsr3_2014}. Para monitoração de exposição externa, o ICRU Report 95 propõe a revisão das grandezas operacionais utilizadas em proteção radiológica\supercite{icru95_2020}.

\section{Regulamentação brasileira}

No Brasil, a RDC nº 611/2022 da ANVISA estabelece os requisitos sanitários gerais para organização e funcionamento dos serviços de radiologia diagnóstica e intervencionista e para o controle das exposições decorrentes dessas práticas\supercite{anvisa_rdc611_2022}. A regulamentação é complementada por instruções normativas específicas por modalidade: a IN nº 90/2021 trata da radiografia médica convencional\supercite{anvisa_in90_2021}; a IN nº 91/2021, de fluoroscopia e radiologia intervencionista\supercite{anvisa_in91_2021}; a IN nº 92/2021, de mamografia\supercite{anvisa_in92_2021}; e a IN nº 93/2021, de tomografia computadorizada médica\supercite{anvisa_in93_2021}. As normas subsequentes abrangem radiologia odontológica extraoral\supercite{anvisa_in94_2021}, radiologia odontológica intraoral\supercite{anvisa_in95_2021}, ultrassonografia diagnóstica ou intervencionista\supercite{anvisa_in96_2021} e ressonância magnética\supercite{anvisa_in97_2021}.

Para proteção radiológica de fontes de radiação em sentido mais amplo, a Norma ANSN 3.01 estabelece os requisitos básicos de radioproteção e segurança radiológica\supercite{ansn301_2025}. A Norma ANSN 3.02 disciplina a organização e atuação dos serviços de radioproteção\supercite{ansn302_2026}. Para medicina nuclear, a Norma CNEN NN 3.05 permanece a referência específica para requisitos de segurança e proteção radiológica dos serviços\supercite{cnen305_2013}.


%%% Local Variables:
%%% mode: LaTeX
%%% TeX-master: "../main"
%%% End:
